\pdfbookmark[1]{Resumo}{resumo}
\chapter*{Resumo}

Sistemas biológicos exibem organização multi-escala onde vias metabólicas
executam transformações bioquímicas enquanto redes regulatórias controlam a
ativação dessas vias. As abordagens computacionais existentes modelam essas
escalas separadamente, impedindo a predição de limiares de decisão ---
fronteiras quantitativas onde células transitam irreversivelmente entre
fenótipos. Este trabalho apresenta a Teoria das Redes de Petri Hierárquicas de
Sinais, uma extensão formal das Redes de Petri clássicas que possibilita
a modelagem unificada de metabolismo e regulação mediante semântica
composicional.

O formalismo estende a 5-tupla clássica $(P, T, F, W, M_0)$ a uma 13-tupla
que incorpora quatro inovações essenciais.
Os \emph{lugares de sinal} ($\Psi \subset P$) particionam os metabólitos entre
aqueles que executam transferência de massa e os que intermediam controle
regulatório, permitindo que moléculas como o ATP participem simultaneamente da
glicólise (por arcos normais) e do controle do desenvolvimento (por arcos de
fluxo de sinal).
Os \emph{arcos de fluxo de sinal} ($F_s$) implementam transformação de
informação com estequiometria, criando estados regulatórios em cascata através
das camadas hierárquicas.
As \emph{funções de camada hierárquicas}
($\lambda: \Psi \to \mathbb{N}$) atribuem níveis organizacionais com regras de
precedência que garantem que camadas inferiores (metabolismo) preemptam camadas
superiores (regulação gênica).
O \emph{parâmetro termodinâmico} $\Gamma = (K, n, \varepsilon)$ determina
analiticamente o limiar efetivo de habilitação por arco
$\theta_{\text{eff}} = K \cdot (\varepsilon / (1 - \varepsilon))^{1/n}$;
o ponto de comprometimento sistêmico --- a concentração e o instante em que a
rede se compromete irreversivelmente --- emerge da dinâmica acoplada da rede
em lugar de ser ajustado como parâmetro.

A validação do formalismo por meio da esporulação de \emph{Bacillus subtilis},
bem caracterizada na literatura de biologia de sistemas, demonstra capacidade
preditiva quantitativa. A construção de um modelo com 40~lugares e
36~transições a partir de cinéticas enzimáticas publicadas (banco de dados BRENDA)
e estequiometria de fosforretransmissão (literatura bioquímica) reproduziu a
assimetria experimental de \textcite{Fujita2005}: a acumulação gradual de
Spo0A${\sim}$P via fosforretransmissão natural produz esporulação em
${\approx}\,52\%$ das réplicas, enquanto a indução constitutiva abrupta de
Spo0A* (Spo0A-Sad67) produz apenas ${\approx}\,5\%$ --- a separatriz de
comprometimento emerge quando $\sigma^H$ ultrapassa $\theta_{\sigma H}
\approx 1{,}60\,\mu$M (varredura fatorial $2 \times 3 \times 4$, 25~condições,
1250~simulações estocásticas) --- sem ajuste de qualquer parâmetro ao
resultado experimental.
A simulação gerou todas as 16~espécies regulatórias e estruturais da cascata
de esporulação, incluindo esporo maduro, com perfis de concentração e
temporização biologicamente plausíveis.

O formalismo aborda limitações da biologia de sistemas computacional além da
decisão bacteriana. Dinâmicas heterogêneas --- cinética enzimática contínua,
expressão gênica estocástica, pontos de verificação temporais --- executam-se
dentro de uma semântica unificada. A arquitetura regulatória torna-se
topologicamente explícita por meio dos tipos de arco, em vez de implícita nas
equações. A construção automatizada de modelos integra bancos de dados de vias
(KEGG) e enzimas (BRENDA), reduzindo o esforço manual e preservando a exatidão
formal mediante verificação de equilíbrio elementar.

Metodologicamente, as Redes de Petri Hierárquicas de Sinais estabelecem que
semânticas composicionais formais podem abranger escalas organizacionais,
possibilitando uma biologia de sistemas integrada onde limiares de
decisão celular emergem da arquitetura da rede.

\vspace{1cm}
\noindent\textbf{Palavras-chave:} Redes de Petri Hierárquicas de Sinais.
Vias Bioquímicas. Redes de Regulação Gênica. Redes de Petri Bio-Híbridas.
Biologia de Sistemas. Biologia Sintética. Modelagem \textit{in~silico}.

\cleardoublepage

% ── Abstract (inglês) ────────────────────────────────────────
\pdfbookmark[1]{Abstract}{abstract}
\begin{otherlanguage}{english}
\chapter*{Abstract}

Biological systems exhibit multi-scale organization where metabolic pathways
execute biochemical transformations while regulatory networks control pathway
activation. Existing computational approaches model these scales separately,
preventing prediction of commitment thresholds---quantitative boundaries where
cells transition irreversibly between phenotypes. This thesis introduces Signal
Hierarchical Petri Net Theory, a formal extension of classical Petri nets
enabling unified modeling of metabolism and regulation through compositional
semantics.

The formalism extends the classical 5-tuple $(P, T, F, W, M_0)$ to a 13-tuple
incorporating four key innovations: signal places ($\Psi \subset P$) partitioning
metabolites into mass-transfer versus regulatory-control roles; signal flow arcs
($F_s$) implementing consumptive information transformation across hierarchical
layers; hierarchical layer functions ($\lambda: \Psi \to \mathbb{N}$) with
precedence rules enforcing that lower layers (metabolism) preempt higher layers
(gene regulation); and the thermodynamic parameter
$\Gamma = (K, n, \varepsilon)$ determining the effective enablement threshold
$\theta_{\text{eff}} = K \cdot (\varepsilon / (1 - \varepsilon))^{1/n}$
that emerges from network dynamics rather than being fitted.

Validation using \emph{Bacillus subtilis} sporulation demonstrates quantitative
predictive capacity. A model with 40~places and 36~transitions, parametrized
exclusively from published enzymatic kinetics (BRENDA) and phosphorelay
stoichiometry, reproduces the experimental asymmetry reported by
\textcite{Fujita2005}: gradual Spo0A${\sim}$P accumulation via natural
phosphorelay produces sporulation in ${\approx}\,52\%$ of replicates, while
constitutively active Spo0A* (Spo0A-Sad67) produces only ${\approx}\,5\%$ ---
the commitment separatrix emerges when $\sigma^H$ exceeds
$\theta_{\sigma H} \approx 1.60\,\mu$M ($2 \times 3 \times 4$ factorial sweep,
25~conditions, 1250~stochastic simulations) --- without fitting any parameter
to the experimental result. The simulation generates all 16~regulatory and
structural species of the sporulation cascade, including mature spore, with
biologically plausible concentration profiles and timing.

\vspace{1cm}
\noindent\textbf{Keywords:} Signal Hierarchical Petri Nets. Biochemical Pathways.
Gene Regulatory Networks. Bio-Hybrid Petri Nets. Systems Biology. Synthetic
Biology. In-Silico Modeling.
\end{otherlanguage}
